\section{Finite Pearson functionals and uniform nonvanishing}\label{sec:rank}
\subsection{Classification of finite solutions}
A linear functional on $\C[z]$ is called finite if it factors through
$\C[z]/(p)$ for some nonzero polynomial $p$.  Equivalently, it is a finite
sum of evaluation derivatives at finitely many points.  Its local component
at a point is a jet; multiplication by a polynomial nonzero at that point
is invertible on the local jet space.  Multiplication by a zero of order
$d$ lowers a nonzero highest derivative order $h$ to $h-d$, or kills
the jet if $h<d$.

\begin{proposition}\label{prop:finite-classification}
For the endpoints~\eqref{eq:endpoints}, the space of finite functionals $U$
satisfying
\begin{equation}\label{eq:finite-pearson}
 U(f(-z))=-U(f(z)),\qquad
 U\bigl(A(z)q((z+1)^2)\bigr)=0\quad(q\in\C[X])
\end{equation}
has dimension seven.  Each solution is supported on the even integers
between $-134n$ and $134n$ and is determined by its odd center jets of
orders $1,3,\ldots,13$.
\end{proposition}
\begin{proof}
The polynomials $q((z+1)^2)$ are precisely the polynomials invariant under
$\tau(z)=-2-z$.  Consequently the second identity says that $AU$ is odd
under $\tau$.  The first says that $U$ is odd under $\sigma(z)=-z$.
Away from a zero of $A$, these identities propagate a support point by
$\sigma\tau(z)=z+2$.  A finite nonempty chain must start at $-a$ and end
at $a'$ for roots $a,a'$ of $A$, with $a+a'\ge0$.

The roots of $A$ are $-200n-2$ and the positive even integers $2a_i$.
A pair involving the negative root has negative sum, since $2a_i\le134n$.
Thus both endpoints arise from positive roots.  Every chain lies on the
even integer lattice, in the indicated interval, and crosses zero.
There can be no finite cycle under translation by two over $\C$.

On the nonnegative half-lattice, the local equation at consecutive points
can be written as equality of the translated jets obtained by multiplying
at $2j$ by $A(z)$ and at $2j+2$ by $A(-z)$.  The latter factor is a unit
because $A(-2j-2)\ne0$ for $0\le j\le67n$.  Thus a center jet determines
every later jet uniquely.  Its highest order decreases by the multiplicity
of each positive root crossed.  There are fifteen such roots counted with
multiplicity.  A highest center order below fifteen is killed by the
terminal endpoint; an order at least fifteen survives and cannot yield
finite support.  Oddness allows only odd orders, giving exactly
$1,3,\ldots,13$.  Conversely, each such center jet propagates and
terminates; reflection provides the negative half and satisfies the same
local equations.  This proves both the dimension and the claimed
parametrization.  In particular, the five coincident roots at $134n$
lower the order by five; distinctness is not assumed.
\end{proof}

\begin{proposition}\label{prop:finite-restrictions}
For $\ell_U(P)=U(zP(z^2))$, the restrictions of all solutions in
Proposition~\ref{prop:finite-classification} to polynomials of degree at
most eight are exactly
\begin{equation}\label{eq:finite-span}
 \operatorname{span}_{\C}\{B_3,B_5,B_7,B_9,C_1,C_3\}.
\end{equation}
Here and below the fixed change $X=4Y$ is understood.  This subspace has
dimension six and contains no functional with nonzero $C_4$ coordinate
in the basis~\eqref{eq:nine}.
\end{proposition}
\begin{proof}
The even rational function $g(z)=w_n(z/2)/z$ satisfies
\[
 g(z+2)/g(z)=A(z)/A(-z-2)
\]
and has a center pole of order ten.  At every even integer use the same
even Laurent germ in the local coordinate.  Internal residues of $g$
times any such germ satisfy the local Pearson propagation equation.
The germs $1,u^2,u^4,u^6,u^8$ give, after multiplication of the test
polynomial by $z$, nonzero scalar multiples of $B_1,B_3,B_5,B_7,B_9$.
Their highest center derivative orders are $9,7,5,3,1$.
Nonnegative-power germs are merely local data; a globally defined
periodic function with these exact germs is not required.

Two further choices are the globally periodic kernels
$(\pi/2)^2\csc^2(\pi z/2)$ and $(\pi/2)^4\csc^4(\pi z/2)$.
Their highest center derivative orders are eleven and thirteen.  Hence
the seven residue functionals just described form a basis of the finite
solution space, by their distinct highest center orders.

It remains to identify the two additional restrictions.  This is most
transparent in the physical $x$ coordinate.  Put
\[
 K_2(x)=\pi^2\csc^2(\pi x),\qquad
 K_4(x)=\pi^4\csc^4(\pi x).
\]
Their principal parts at an integer are respectively $u^{-2}$ and
$u^{-4}+(2\pi^2/3)u^{-2}$.  Integrate
$w_n(x)P(x^2)K_h(x)$ on expanding rectangles with vertical sides
at half-integers.  For $\deg P\le8$, rational decay makes the vertical
integrals tend to zero, while the horizontal integrals decay
exponentially.  The sum of all residues is therefore zero.  The zeros
of order five at $a_0<|m|\le b$ remove every intervening residue.
At the remaining external integers reflection doubles the positive
half of the sum.  Thus the two internal residue functionals, in this
coordinate, are
\begin{equation}\label{eq:kernel-residues}
 -2F_1(P),\qquad -2F_3(P)-\frac{4\pi^2}{3}F_1(P).
\end{equation}
Internal residues must include every integer $|m|\le a_0$, even when
$w_n$ is regular there: the trigonometric kernel still has a pole.
In particular the nodes $63n<|m|\le67n$ cannot be omitted.

The partial-fraction formula expresses $F_1$ and $F_3$ as $C_1$ and
$C_3$, respectively, plus linear combinations of the four nonzero $B$
directions.  Only even zeta values occur in these combinations, and
their values are irrelevant to the span.  Also $B_1=0$ in the proper
range.  Equation~\eqref{eq:kernel-residues} is triangular with nonzero
coefficients on $C_1,C_3$.  Lemma~\ref{lem:independent} now proves the
dimension and the final exclusion assertion.
\end{proof}

\subsection{A truncated Hankel extension lemma}
Hankel rank is closely related to binary apolarity. The following standard
consequence is included with
its degree bookkeeping, since a possible root at infinity matters here.
\begin{lemma}\label{lem:hankel}
Let $\ell$ be a nonzero linear functional on $\C[X]_{\le2m}$, and
suppose its middle Hankel matrix $(\ell(X^{i+j}))_{0\le i,j\le m}$
has rank $r\le m$.  There is a nonzero polynomial $p$ of degree
$s\le r$ and a functional $\widetilde\ell$ on $\C[X]$, annihilated
by $p$, which agrees with $\ell$ through degree $2m-r+s$.
For $s=0$ this means agreement with the zero functional.
\end{lemma}
\begin{proof}
Write $m_j=\ell(X^j)$ and homogenize the moments to the binary form
\[
 f(U,V)=\sum_{j=0}^{2m}\binom{2m}{j}m_jU^{2m-j}V^j.
\]
Its catalecticant maps have Hankel matrices in the divided-power
normalization.  The apolar ideal of a binary form is a complete
intersection with generator degrees $d,2m+2-d$, where $d\le m+1$;
see \cite[Proposition 1.7]{BS}.  Its Hilbert function in degree $m$
is $d$, so $d=r$ when the middle rank is at most $m$.  The case where
$f$ is a pure power has the same conclusion directly, with $r=1$.

A degree-$r$ homogeneous annihilator, dehomogenized, is a polynomial
$p(X)=\sum_{j=0}^sp_jX^j$ with $s\le r$ and $p_s\ne0$.  Its
catalecticant equations are
\[
 \sum_{j=0}^sp_jm_{j+k}=0\qquad(0\le k\le2m-r).
\]
Extend the first $s$ moments by this recurrence.  The resulting
functional is annihilated by $p$, and the displayed equations force
agreement through degree $2m-r+s$.  If $s=0$, those equations say
directly that all moments in the asserted range are zero.  The drop
$r-s$ is exactly what permits homogeneous factors at infinity; no
flat-extension assumption has been inserted. The Lean development proves
the needed short-annihilator and recurrence-extension statements directly;
it does not assume an apolarity theorem as an external axiom.
\end{proof}

\begin{theorem}\label{thm:rank}
Let $m\ge4$ and $2m\le\beta+3$.  For arbitrary complex numbers
$c_3,c_5,c_7,c_9$, put $\ell=C_4+\sum c_jB_j$.  Then
\[
 \rank(\ell(Y^{i+j}))_{0\le i,j\le m}\ge(m+1)-4.
\]
\end{theorem}
\begin{proof}
Assume the contrary, so its rank $r$ is at most $m-4$.
Apply Lemma~\ref{lem:hankel}, and put $e=r-s$.  The extension agrees
through degree $2m-e\ge m+4\ge8$.  It cannot be zero by
Lemma~\ref{lem:independent}, so $s>0$.
Define a functional on $\C[z]$ by
\[
 U(f)=\widetilde\ell\left(\frac{f(z)-f(-z)}{2z}\right),
\]
reading the even quotient as a polynomial in $X=z^2$.  It is odd and
annihilated by $p(z^2)$, and therefore factors through an algebra of
dimension $2s$.  Let $W(q)=U(A(z)q((z+1)^2))$.
By definition, $W(q)=\widetilde\ell(\mathcal Dq)/2$.
For $0\le k\le2s-1$,
\[
 \deg\mathcal D(X^k)\le2s+8\le2m-(r-s),
\]
since $2m-r-s\ge2(m-r)\ge8$.  These moments vanish by agreement
with $\ell$ and Lemma~\ref{lem:pearson}.  Multiplication by
$(z+1)^2$ in $\C[z]/(p(z^2))$ is an endomorphism of a
$2s$-dimensional space.  Cayley--Hamilton extends the vanishing to
every $k\ge0$.  Hence $U$ is a finite Pearson solution.  Its first
nine induced moments agree with $\ell$, contradicting
Proposition~\ref{prop:finite-restrictions}, since the $C_4$ coordinate
of $\ell$ is one.
\end{proof}

\subsection{Selection of a nonzero minor}
Put
\begin{equation}\label{eq:N-K}
 M=100n,\qquad C=\frac{361}{200},\qquad
 N=\floor{361n/2}+1,\qquad K=N-4.
\end{equation}
Then $2N-2\le\beta$ and $(N-1)/M\le C$.  We leave the final
few permissible polynomial degrees unused to retain the latter exact
inequality.  Define
\begin{equation}\label{eq:E}
 E_0(Y)=1,\qquad
 E_i(Y)=\frac{2\prod_{h=0}^{i-1}(Y-h^2)}{(2i)!}\quad(i\ge1).
\end{equation}
These form a basis in every initial degree range.  Moreover
\begin{equation}\label{eq:E-integer}
 E_i(x^2)=\binom{x+i}{2i}+\binom{x+i-1}{2i}\quad(i\ge1),
\end{equation}
an identity of polynomials, so $E_i(x^2)$ is integer-valued on $\Z$.
One can verify~\eqref{eq:E-integer} by adding the two factorial products:
their common factors give $2x^2\prod_{h=1}^{i-1}(x^2-h^2)/(2i)!$.
Generalized binomial coefficients at negative integer arguments are
integers as well.

Apply Theorem~\ref{thm:rank} to~\eqref{eq:zeta-form}.  The matrix
$(F_4(E_iE_j))_{0\le i,j<N}$ has rank at least $K$.  Hence there
exist sets $I,J\subset\{0,\ldots,N-1\}$ of cardinality $K$ for which
\begin{equation}\label{eq:minor}
 D_n=\det\bigl(F_4(E_iE_j)\bigr)_{i\in I,j\in J}\ne0.
\end{equation}
Under simultaneous rationality it belongs to $\Q$.  The rest of the
argument bounds all these minors uniformly; it imposes no condition
on the chosen $I,J$.
